% Canonical Paper I source.
An arithmetic expression has several mathematically distinct lives.  It is a finite
syntax tree; it may be equipped with a legal order of evaluation; a chosen running
value turns that order into a marked history; the history induces an operator; and
the operator sends an initial value to an endpoint.  Ordinary arithmetic usually
passes through these levels as quickly as possible.  AEG instead asks what is lost
at each passage.

The guiding structure of this paper is
\begin{equation}\label{eq:process-result-hierarchy}
\begin{gathered}
\text{unmarked sequential tree}
  \xrightarrow{\text{choose an accumulator}}
\text{marked planar sequential tree},\\
\text{marked planar sequential tree}
  \longleftrightarrow
\text{bounded marked history},\\
\text{bounded marked history}
  \longrightarrow \text{operator}
  \longrightarrow \text{endpoint value},\\
\text{bounded marked history}
  \longrightarrow \text{charge shadow}.
\end{gathered}
\end{equation}
The first arrow denotes a choice, not a canonical map.  When the planar structure,
mark, operation labels, constants, and slots are retained, the double arrow is the
bijection proved in Section~\ref{sec:sequential-histories}.  The two arrows leaving a
history are separate forgetful maps: charge is not a function of the evaluated
operator, and an endpoint is not a function of charge alone.  Different histories
may nevertheless induce the same operator, and different operators may agree at one
chosen initial value.  The geometry belongs primarily to these intermediate
levels, before endpoint evaluation has erased the order of construction.

Figure~\ref{fig:process-result-structure} renders this structure explicitly.
It is a map of the distinctions used in the paper, not a claim that any
forgetful arrow is injective.

\begin{figure}[t]
\centering
\begin{tikzpicture}[
  stage/.style={draw=black!60,rounded corners=3pt,fill=blue!5,
                minimum height=9mm,align=center,font=\small,inner xsep=5pt},
  forget/.style={-{Latex[length=2mm]},thick,black!65},
  choice/.style={-{Latex[length=2mm]},thick,dashed,orange!70!black},
  bij/.style={{Latex[length=2mm]}-{Latex[length=2mm]},thick,blue!60!black}
]
  \node[stage] (tree) at (-4.5,1.25) {unmarked\\sequential tree};
  \node[stage] (marked) at (-1.0,1.25) {marked planar\\sequential tree};
  \node[stage,fill=blue!10] (history) at (2.6,1.25)
    {bounded marked\\history};
  \node[stage,fill=violet!8] (charge) at (0,-0.8) {charge shadow};
  \node[stage] (operator) at (3.35,-0.8) {operator};
  \node[stage] (endpoint) at (5.75,-0.8) {endpoint\\value};

  \draw[choice] (tree)--node[above=5pt,font=\scriptsize] {choose accumulator} (marked);
  \draw[bij] (marked)--node[above=5pt,font=\scriptsize] {encode/reconstruct} (history);
  \draw[forget] (history)--node[pos=0.58,right=3pt,font=\scriptsize]
    {evaluate word} (operator);
  \draw[forget] (operator)--(endpoint);
  \draw[forget] (history)--node[pos=0.55,left=4pt,font=\scriptsize]
    {record charges} (charge);
\end{tikzpicture}
\caption{The process-to-result structure.  Choosing an accumulator marks an
unmarked sequential tree; a fully marked planar tree and a bounded marked
history determine one another.  From the history, operator evaluation and the
charge map are separate information-losing branches, while an endpoint is
obtained by applying the operator to an initial value.}
\label{fig:process-result-structure}
\end{figure}

\subsection{Sequential histories}

Not every expression tree determines a unique history.  Independent subexpressions
may be evaluated in either order, so a single tree can have many legal linear
extensions of its dependency order.  Section~\ref{sec:sequential-histories} isolates
the precise exception.  For a finite binary tree, uniqueness of the internal
evaluation order is equivalent to the internal vertices forming one chain, or
equivalently one spine.  A marked accumulator on that spine determines a word of
one-hole contexts
\[
 C_{\omega,c}^{(1)}[z]=\omega(z,c),
 \qquad
 C_{\omega,c}^{(2)}[z]=\omega(c,z).
\]
The superscript records the operand slot occupied by the current value.  It is not a
time direction.  This distinction prevents three operations that are often blurred
in informal tree pictures---planar mirror, temporal reversal, and path inverse---from
being identified.

This use of syntax is modest.  Formal generation and rewriting have long histories
in logic and theoretical computer science; classical points of reference include
Post's work on combinatorial reduction and Chomsky's generative hierarchies
\cite{Post1943FormalRO,Chomsky1956ThreeMF}.  Here the purpose is different: the
sequential fragment supplies a controlled domain on which evaluation can be treated
as a path without pretending that every arithmetic tree has a unique path.

\subsection{Why bilateral completion matters}

Earlier affine calculations use the current value in one distinguished operand slot.
That fragment is closed under translations and nonzero scalings and is therefore
naturally represented by the affine group of the line.  Once both operand slots are
allowed, the context $z\mapsto c/z$ introduces inversion.  The correct algebraic
ambient space is then the projective line and its fractional linear transformations.

Section~\ref{sec:projective-affine} proves, over an arbitrary field $K$, that the
non-degenerate bilateral contexts generate $\PGL_2(K)$.  The affine transformations
are exactly the stabilizer of the projective point at infinity.  Thus the continuous
geometry is neither discarded nor promoted to the whole theory: it is identified
precisely as the affine sector of a larger projective semantics.  The geometric
reading of these operators --- addition and multiplication as motions of an
arithmetic expression space, and second-slot division as the projective step --- is
developed in Paper~0 and used here only at the operator level.

Inside that larger semantics, the two contexts
\[
  T(z)=z+\sqrt2,
  \qquad
  J(z)=-\frac1z
\]
generate the $q=4$ Hecke group.  Section~\ref{sec:projective-affine} verifies
the relation $(JT)^4=1$ directly from the context matrices and formulates the
resulting interface from arithmetic words to group operators and then to cell
cosets.  This restricted alphabet is important because its hyperbolic
cellulations are produced by arithmetic operator semantics rather than imposed
after endpoint evaluation.  The interface is deliberately not a bijection:
relations identify different histories, and stabilizers identify different
operators at the level of cells or endpoints.

Projective continuation does not repair an inadmissible ordinary computation.  For
example, $c/z$ has a well-defined projective value at $z=0$, namely infinity, while
ordinary division by zero remains undefined.  We keep the ordinary domain and the
projective extension as separate data throughout.

\subsection{Main results}

The first group of results is combinatorial and syntactic.  The trees with a
unique legal internal evaluation order are exactly those whose internal vertices
form a single spine; marking a leaf at the innermost vertex turns such a tree into
a bounded marked spinal history, and the two constructions are mutually inverse
when planar data are retained.  Operand-slot chirality is part of that data and is
kept distinct from chronological direction.

The second group is projective.  The non-degenerate bilateral contexts generate
$\PGL_2(K)$ for every field $K$, the contexts fixing infinity form the affine
sector, and the proof yields the rank-one Bruhat decomposition.  The
$\mathbb Q(\sqrt2)$-sublanguage generated by translation by $\sqrt2$ and by
inversion carries the order-four relation $(JT)^4=1$, realized by an explicit
eight-step marked history that is projectively trivial and ordinarily inadmissible
at one initial value.  That example is the reason this paper keeps ordinary
admissibility as a separate hypothesis from projective non-degeneracy.

The third group is the zero theory of regular arithmetic expression spaces.
Section~\ref{sec:regular-aes-interface} restates the definition imported from
Paper~0; Section~\ref{sec:zero-geometry} proves that the zero locus of a regular
assignment is a smooth curve, that a complete connected boundaryless regular AES
splits as $Z(a)\times\mathbb R$, that this forbids branching or multiple zero
components in the regular category, and that a smooth parameter family with
spatially regular zeros has a submersive total zero set --- with the properness
warning that prevents the last statement from being read as triviality.  The
section then defines singular arithmetic expression spaces and verifies one
isolated-zero model in detail.

Everything else that earlier versions of this paper developed --- the two affine
translation coordinates, the continuous arithmetic motion, the basic hyperbolic
model $\mathfrak E_0$, the accumulated commutative space, tearing, and the
contact structure --- now belongs to Paper~0, where it is proved at the same
labels.  This paper cites those results rather than restating them, with the sole
exception of the regular-AES interface needed to state its own zero theorems.

\subsection{Scope and relation to standard geometry}

Several ingredients used here are standard: finite posets and their linear
extensions, fractional linear transformations, and the regular-value theorem.
The mathematical contribution sought in Paper~I is their history-sensitive
arithmetic organization and the exact interfaces among them.

Four boundaries are enforced.  First, the geometric construction of arithmetic
expression spaces, their motions, their commutativization, and their contact
structure belongs to Paper~0.  Second, a horizontal contact distribution does not
by itself select a unique complex structure, so arithmetic holomorphic and
harmonic theories are deferred to Paper II; this includes any automorphic-function
theory built on the $G_4$ action.  Third, the regular-value theorem is only the
starting point for singular families: pulling analytic real loci back to the Hecke
cellulation, and developing their multiple zeros, discriminants, tubes, braids,
and knots, belong to Paper III.  Fourth, projective bivaluations, homogeneous
quotient towers, and condensation belong to Paper IV, and no computational lower
bound follows from noncommutativity or negative curvature alone.

\subsection{Organization}

Section~\ref{sec:sequential-histories} develops sequential trees and marked
histories.  Section~\ref{sec:projective-affine} gives the bilateral projective
semantics and locates the affine sector.
Section~\ref{sec:regular-aes-interface} restates the regular-AES interface
imported from Paper~0.  Section~\ref{sec:zero-geometry} develops the regular and
singular zero geometry.  Section~\ref{sec:conclusion} records what is proved,
what remains open, and the interfaces to the later papers.  The appendices
collect convention checks and the small examples that separate the equality
levels of \eqref{eq:process-result-hierarchy}.
